\section{Moteur 1 --- Mean Reversion Intraday Filtré (80\,\%)}
\label{sec:mr_macro}

\lettrine[lines=2, findent=2pt, nindent=0pt]{L}{e} premier moteur est la composante dominante du portefeuille~: il reçoit 80\,\% de l'allocation statique. Ce poids prépondérant est justifié par le fait qu'il possède simultanément le Sharpe le plus élevé (environ 0.82 full-period) et la volatilité intrinsèque la plus basse (environ 1.3\,\% annualisée) parmi les trois moteurs. Les deux autres composantes jouent un rôle de diversificateurs à 10\,\% chacune, et leur valeur apparaît surtout dans les années où ce premier moteur est dormant.

\subsection{Intuition économique}

\begin{thesisbox}
\textbf{La \emph{mean reversion} intraday fonctionne dans les régimes de liquidité normale et de transmission monétaire saine.} Lorsque la courbe des taux s'inverse ou que le taux de chômage commence à remonter, les micro-structures qui rendent la \emph{mean reversion} profitable disparaissent~: les flux deviennent unidirectionnels en raison de la \emph{fuite vers la qualité}, les spreads s'élargissent, et toute prise de position à contre-courant devient coûteuse. Le filtre macro bloque précisément la stratégie dans ces régimes.
\end{thesisbox}

Techniquement, la stratégie exploite deux asymétries micro-structurelles bien documentées dans la littérature académique sur le marché FX~:

\begin{enumerate}
\item La tendance des prix intraday à revenir vers le \textbf{VWAP} (\emph{Volume-Weighted Average Price}) au cours de la session de chevauchement Londres~+~New York, période durant laquelle les teneurs de marché rebalancent leurs inventaires.
\item La \textbf{surréaction de court terme} aux annonces~: les déviations supérieures à cinq écarts-types autour du VWAP capturent typiquement une réaction excessive à une publication économique ou à un flux institutionnel concentré.
\end{enumerate}

\subsection{Mécanique du signal}

\begin{enumerate}
\item \textbf{VWAP ancré à la journée}~: le calcul est réinitialisé à 00h00 UTC chaque jour, et reprend de façon cumulative.
\item \textbf{Bandes de Bollinger} à 80 périodes et 5 écarts-types, calculées sur la déviation prix moins VWAP. Les bandes ainsi obtenues définissent des seuils dynamiques de retour à la moyenne.
\item \textbf{Filtre session}~: uniquement entre 6h et 14h~UTC, correspondant au chevauchement Londres et New~York où la liquidité atteint son maximum.
\item \textbf{Filtre macro à deux étages}~:
  \begin{itemize}
  \item Spread de courbe des taux 10~ans moins 2~ans inférieur à 0.5~(proxy de stress cyclique)~;
  \item Taux de chômage dont la variation sur trois mois n'est pas haussière (proxy d'absence de fin de cycle d'expansion).
  \end{itemize}
\item \textbf{Règles d'entrée}~: position longue si le prix touche la bande inférieure, position courte si le prix touche la bande supérieure, sous réserve que les filtres session et macro soient tous deux passants.
\item \textbf{Règles de sortie}~: stop-loss à $-0.5\,\%$, take-profit à $+0.6\,\%$, arrêt forcé à 21h~UTC, durée maximale de détention six heures.
\end{enumerate}

\subsection{Performance historique}

\begin{figure}[H]
\centering
\includegraphics[width=\linewidth]{sleeve_mr_macro_equity.png}
\caption{Équité et drawdown du Moteur~1 pris isolément sur la période 2019--2026 (non pondéré, non leveragé). La ligne verticale pointillée marque le début de la fenêtre hors échantillon.}
\end{figure}

\begin{table}[H]
\centering
\caption{Performance annuelle du Moteur~1 pris isolément. L'année 2021 ne produit aucun trade parce que le spread 10Y--2Y est resté supérieur à 0.5 toute l'année post-Covid, déclenchant le filtre macro.}
\renewcommand{\arraystretch}{1.1}
\begin{tabular}{lrrrr}
\toprule
\textbf{Année} & \textbf{Trades} & \textbf{CAGR} & \textbf{Sharpe} & \textbf{Max DD} \\
\midrule
2019 & 44 & $-0.82$\,\% & $-0.64$ & $-1.71$\,\% \\
2020 & 3 & $+0.69$\,\% & $+1.19$ & $-0.02$\,\% \\
2021 & \textbf{0} & $0.00$\,\% & $0.00$ & $0.00$\,\% \\
2022 & 30 & $+5.21$\,\% & \textbf{$+2.23$} & $-0.67$\,\% \\
2023 & 19 & $-0.38$\,\% & $-0.28$ & $-1.22$\,\% \\
2024 & 19 & $+2.31$\,\% & \textbf{$+2.08$} & $-0.34$\,\% \\
2025 & 11 & $+0.85$\,\% & $+0.72$ & $-0.71$\,\% \\
\bottomrule
\end{tabular}
\end{table}

\begin{warningbox}
\textbf{Le filtre macro peut réduire la stratégie au silence complet.} En 2021, aucun trade n'a été déclenché parce que la courbe des taux restait exceptionnellement raide après le choc Covid. Ce comportement est voulu~: les tests de sensibilité ont montré que relâcher le seuil à 0.7 ou plus détruit le Sharpe de l'année 2025 (passage de $+0.72$ à $-0.38$). Le \emph{silence sélectif} est une propriété désirable du filtre, pas un bug.
\end{warningbox}

\subsection{Exemples de trades réels}

\begin{table}[H]
\centering
\caption{Exemples de trades MR Macro sur EUR-USD — deux meilleurs, un médian, deux pires. La colonne Durée illustre le design intraday~: la position est dénouée dans la séance.}
\label{tab:trades_mr_macro}
\begin{tabular}{llcrrr}
\toprule
Entrée & Sortie & Sens & Durée & Rendement & P\&L \\
\midrule
2025-03-06 08:17 & 2025-03-06 13:42 & Short & 5.4\,h & 0.57\,\% & \$6,121 \\
2025-02-13 13:41 & 2025-02-13 17:25 & Long & 3.7\,h & 0.57\,\% & \$6,083 \\
2022-08-17 14:04 & 2022-08-17 20:04 & Short & 6.0\,h & 0.10\,\% & \$1,016 \\
2024-12-18 14:01 & 2024-12-18 14:39 & Long & 0.6\,h & -0.53\,\% & \$-5,658 \\
2025-01-20 08:31 & 2025-01-20 08:35 & Short & 0.1\,h & -0.53\,\% & \$-5,673 \\
\bottomrule
\end{tabular}
\end{table}


\begin{examplebox}
\textbf{Anatomie d'un trade gagnant~:} le trade long du 12~février 2025 à 08h31~UTC est entré lorsque le \emph{close} EUR-USD a touché la bande inférieure calculée sur la déviation au VWAP. La session étant ouverte (8h~UTC supérieur à 6h) et le spread 10Y--2Y à $0.32$ (sous le seuil 0.5), les deux filtres étaient passants. Le retour du prix vers le VWAP a déclenché la sortie par take-profit à 11h59~UTC, pour une durée totale de 3~heures~30 et un rendement de $+0.57\,\%$, soit un P\&L de 6\,236~USD sur un capital nominal d'un million.
\end{examplebox}
